Soit $ABC$ un triangle. Le cercle exinscrit au triangle $ABC$ opposé au sommet $A$ est tangent au côté $BC$ au
point $A_1$. Les points $B_1$ sur $CA$ et $C_1$ sur $AB$ sont définis de la même façon, en utilisant les
cercles exinscrits opposés à $B$ et $C$ respectivement. On suppose que le centre du cercle circonscrit au
triangle $A_1B_1C_1$ est un point du cercle circonscrit au triangle $ABC$. Montrer que le triangle $ABC$ est
rectangle.

\emph{Le cercle exinscrit au triangle $ABC$ opposé au sommet $A$ est le cercle tangent au segment $[BC]$, à la
demi-droite $[AB]$ au-delà de $B$ et à la demi-droite $[AC]$ au-delà de $C$. Les cercles exinscrits opposés à
$B$ et à $C$ sont définis de la même manière.}
